\chapter{Offspring laws and extinction}
\label{ch:offspring}

\begin{definition}[Offspring law]
\label{def:offspring}
\lean{BranchingProcess.Offspring, BranchingProcess.Offspring.mean}
\leanok
An \emph{offspring law} with support bound $J$ is a sequence $\theta=(\theta_j)_{j\geq0}$ of
nonnegative reals vanishing above $J$ and summing to one. Its \emph{mean} is
$m=\sum_j j\,\theta_j$, and $\theta$ is \emph{subcritical}, \emph{critical} or
\emph{supercritical} according as $m<1$, $m=1$ or $m>1$.
\end{definition}

\begin{definition}[Generating function]
\label{def:gen}
\uses{def:offspring}
\lean{BranchingProcess.Offspring.gen, BranchingProcess.Offspring.genDeriv}
\leanok
The generating function of $\theta$ is $f(s)=\sum_{j=0}^{J}\theta_j s^j$.
\end{definition}

\begin{lemma}[Elementary properties of $f$]
\label{lem:gen}
\uses{def:gen}
\lean{BranchingProcess.Offspring.gen_one, BranchingProcess.Offspring.gen_le_one,
BranchingProcess.Offspring.gen_mono, BranchingProcess.Offspring.hasDerivAt_gen,
BranchingProcess.Offspring.genDeriv_one}
\leanok
$f(1)=1$, $f(0)=\theta_0$, $f$ maps $[0,1]$ into itself monotonically, and $f'(1)=m$.
\end{lemma}

\begin{proof}
\leanok
Each is a termwise computation on the finite sum.
\end{proof}

\begin{definition}[Extinction probability]
\label{def:extinction}
\uses{def:gen}
\lean{BranchingProcess.Offspring.extinction}
\leanok
The \emph{extinction probability} of $\theta$ is
\[
  q=\inf\set{s\in[0,1] : f(s)=s},
\]
the least fixed point of $f$ in $[0,1]$.
\end{definition}

\begin{lemma}[The infimum is a fixed point]
\label{lem:extinction-fixed}
\uses{def:extinction}
\lean{BranchingProcess.Offspring.extinction_mem, BranchingProcess.Offspring.gen_extinction,
BranchingProcess.Offspring.extinction_le_of_fixed}
\leanok
$q\in[0,1]$, $f(q)=q$, and $q\leq s$ for every fixed point $s\in[0,1]$ of $f$.
\end{lemma}

\begin{proof}
\leanok
The fixed-point set is closed, nonempty because $f(1)=1$, and bounded below by $0$, so it
contains its infimum. Minimality is the defining property of the infimum.
\end{proof}

\begin{lemma}[Extinction and childlessness]
\label{lem:extinction-zero}
\uses{lem:extinction-fixed}
\lean{BranchingProcess.Offspring.extinction_eq_zero_iff, BranchingProcess.Offspring.extinction_pos}
\leanok
$q=0$ if and only if $\theta_0=0$; in particular $\theta_0>0$ forces $q>0$.
\end{lemma}

\begin{proof}
\leanok
If $q=0$ then $\theta_0=f(0)=f(q)=q=0$. Conversely $\theta_0=0$ makes $0$ a fixed point, and
minimality gives $q\leq0$.
\end{proof}

\begin{theorem}[A supercritical law survives]
\label{thm:supercritical}
\uses{lem:extinction-fixed, lem:gen}
\lean{BranchingProcess.Offspring.extinction_lt_one_of_supercritical}
\leanok
If $m>1$ then $q<1$.
\end{theorem}

\begin{proof}
\leanok
Write $1-f(s)=(1-s)g(s)$ with $g(s)=\sum_j\theta_j(1+s+\dots+s^{j-1})$, a polynomial with
$g(1)=m>1$. By continuity $g>1$ on a left neighbourhood of $1$, so $f(s)<s$ for some $s<1$.
Since $f(0)=\theta_0\geq0$, the intermediate value theorem produces a fixed point in $[0,s]$,
and minimality pushes $q$ below it.
\end{proof}
